\documentclass[10pt,a4paper]{amsart}
\usepackage[margin=19mm]{geometry}
\usepackage{amsmath,amssymb,mathtools}
\usepackage[scr=rsfs,cal=euler]{mathalfa}
\usepackage{xfrac}
\usepackage[unicode,hidelinks,bookmarks=false]{hyperref}
\newcommand{\ie}{i.e.\ }
\newcommand{\cf}{cf.\ }
\newcommand{\resp}{respectively\ }
\newcommand{\Subsection}{\subsection}
\providecommand{\nobreakdash}{\nobreak\hbox{-}}
\setlength{\emergencystretch}{2em}
\multlinegap0pt
\usepackage{bm}
\usepackage{bbm}
\usepackage{dsfont}
\usepackage{xspace}
\usepackage{cancel}
\usepackage{mathtools}
\usepackage{amsthm}
\makeatletter
\@ifundefined{definition}{\newtheorem{definition}{Definition}}{}
\@ifundefined{theorem}{\newtheorem{theorem}{Theorem}}{}
\makeatother
\title[A generalization of Deterministic Finite Automata related to]{A generalization of Deterministic Finite Automata related to discharging}
\author{John M. Campbell}
\date{Source: arXiv:2506.14072v1 (CC BY 4.0)}
\begin{document}
\maketitle

\noindent\emph{Source and licence.} A generalization of Deterministic Finite Automata related to discharging. Version arXiv:2506.14072v1 (CC BY 4.0), distributed under the Creative Commons Attribution 4.0 International licence, as recorded in the arXiv licence metadata for this exact version and shown on the abstract page https://arxiv.org/abs/2506.14072v1. Full rights notice: licenses/arxiv-2506.14072-CC-BY-4.0.txt.\par\smallskip

\noindent\textit{Editorial scope.} Counted unit: the Theorem whose source label is \texttt{theoremmainformalize} in the article (source0.tex lines 756--764, source label \texttt{theoremmainformalize}), delivered together with the article's own complete original proof of it (source0.tex lines 766--857, one \texttt{proof} environment of 92 lines). No mathematics is written, repaired or re-derived: statement and proof are the article's own text line for line. Required context retained uncounted: definitionchargeext (lines 430--488); deltanested (lines 446--448); setdeltacqw (lines 480--484); reduceddefinition (lines 726--742). Every definition, display and label of those lines is retained unchanged. Omitted from this package and not redistributed: 1--429, 449--479, 485--725, 743--755, 765--765, 858--1210. Nothing inside a retained excerpt is omitted or altered: every retained body is delivered exactly as the article's lines are written, so no omission receipt of any kind accompanies this package. Citations: the retained excerpts carry no citation command at all, so no bibliography entry of the article is transcribed and no reference is redistributed. Reference adaptation: none was needed and none was made; every reference inside the delivered unit resolves inside it, so no unresolved reference or citation is produced. Printed slips kept literal: no printed word, symbol, hypothesis or number of the counted statement or of its proof was altered. Layout and class: the article's own class is \texttt{article} with packages including amssymb, amsmath, amsthm, url, tikz; the wrapper is the lane's pinned \texttt{amsart} editorial wrapper at 10pt on A4, which restates the theorem-like environments the retained text needs and adds the article's own macro definitions that the retained text needs (none), with extra wrapper packages (bm, bbm, dsfont, xspace, cancel, mathtools). The delivered numbering is the unit's own. Assets: no retained span contains a figure environment, a graphics-inclusion command, a PDF-page inclusion, a table or a TikZ picture; no image file is copied into this package, the package declares no asset dependency and no image file is redistributed with it. Licence: version arXiv:2506.14072v1 of this article is distributed under the Creative Commons Attribution 4.0 International licence, as recorded in the arXiv licence metadata for this exact version and shown on the abstract page https://arxiv.org/abs/2506.14072v1. A full rights notice is delivered at licenses/arxiv-2506.14072-CC-BY-4.0.txt. Characters: every retained excerpt is pure ASCII, so the delivered engine, XeLaTeX with its default Latin Modern text font, reports no missing character.
\begin{definition}\label{definitionchargeext}
 Letting $Q$ and $\Sigma$ and $\delta$ and $q_0$ and $P$
 be as in the definition for a DDFA, 
 we define the 
 \emph{charge-extended transition function}
 $$ \delta^{c}\colon Q \times \Sigma^{\ast} \to Q \times [0, 1] $$
 as follows. For distinguished $q \in Q$, we set $\delta^{c}(q, \epsilon) = (q, 1)$. 
 For a nonempty word $w \in \Sigma^{\ast}$ we write $w = w_{1} w_{2} \cdots w_{\ell(w)}$. 
 We also let $Q = \{ q^{(1)}$, $ q^{(2)}$, $ \ldots$, $ q^{(|Q|)} \}$, 
 and we define the \emph{contemporary charge} $c^{q, i}$ of an element 
 in $Q$ and parameterized by $q$ 
 and $i \in \{ 0, 1, 2, \ldots, \ell(w) \}$ as follows. 
 To begin with, we set $c^{q, 0}(q) = 1$ 
 and $c^{q, 0}(q^{(j)}) = 0$ for $q^{(j)} \neq q$. 
 For $i \in \{ 1, 2, 3, \ldots, \ell(w) \}$, we then let $k_i$ denote the unique 
 integer in $\{ 1, 2, \ldots, |Q| \}$ such that 
\begin{equation}\label{deltanested}
 q^{(k_i)} = \delta(\ldots \delta(\delta(\delta(q, w_1), w_2), w_3), \ldots, w_{i}), 
\end{equation}
 and we also write $q^{(k_0)} = q$. 
 Adopting the convention whereby a sum over the null set is $0$, we set 
\begin{multline*}
 c^{q, i+1}\big( q^{(k_{i})} \big) = \\
 \begin{cases} 
 c^{q, i}\big( q^{(k_{i})} \big) \left( 1 - n^{\text{current}}_{q^{(k_i)}, w_{i+1}} 
 - \sum_{ \substack{ t \neq w_{i+1} \\ \delta(q^{(k_i)}, t) \neq q^{(k_i)} } } n_{q^{(k_i)}, t}^{\text{not current}} \right) & \null \\ 
 \hspace{2.7in} \text{if $\delta\big( q^{(k_i)}, w_{i+1} \big) \neq q^{(k_i)} $}, & \null \\
 c^{q, i}\big( q^{(k_{i})} \big) \left( 1 
 - \sum_{ \substack{ t \neq w_{i+1} \\ \delta(q^{(k_i)}, t) \neq q^{(k_i)} } } n_{q^{(k_i)}, t}^{\text{not current}} \right) & \null \\ 
 \hspace{2.7in} \text{if $\delta\big( q^{(k_i)}, w_{i+1} \big) = q^{(k_i)} $}, & \null 
 \end{cases}
\end{multline*}
 with 
\begin{multline*}
 c^{q, i+1}\big( q^{(k_{i+1})} \big) = \\
 \begin{cases} 
 c^{q, i}\big( q^{(k_{i+1})} \big) + c^{q, i}\big( q^{(k_{i})} \big) \left( n^{\text{current}}_{q^{(k_i)}, w_{i+1}} 
 + \sum_{ \substack{ t \neq w_{i+1} \\ \delta(q^{(k_i)}, t) = q^{(k_{i+1})} } } n_{q^{(k_i)}, t}^{\text{not current}} \right) & \null \\ 
 \hspace{2.7in} \text{if $\delta\big( q^{(k_i)}, w_{i+1} \big) \neq q^{(k_i)} $}, & \null \\ 
 c^{q, i+1}\big( q^{(k_i)} \big) & \null \\ 
 \hspace{2.7in} \text{if $\delta\big( q^{(k_i)}, w_{i+1} \big) = q^{(k_i)} $}, & \null 
 \end{cases} 
\end{multline*}
 and, for $j \in \{ 1, 2, \ldots, |Q| \}$ not equal to $k_i$ and not equal to $k_{i+1}$, we set 
\begin{equation}\label{202506121177A77M1A}
 c^{q, i+1}\big( q^{(j)} \big) = c^{q, i}\big( q^{(j)} \big) 
 + c^{q, i}\big( q^{(k_{i})} \big) \sum_{\substack{t \\ \delta\left( q^{k_i}, t \right) 
 = q^{(j)} }} n_{q^{(k_i)}, t}^{\text{not current}}. 
\end{equation}
 We then set 
\begin{equation}\label{setdeltacqw}
 \delta^{c}( q, w ) 
 = \left( q^{( k_{\ell(w)} )}, 
 c^{q, \ell(w)}\big( q^{( k_{\ell(w)} )} \big) \right). 
\end{equation}
 We also write $ \delta^{c}( q, w, 1 ) = q^{( k_{\ell(w)} )}$ and 
 $ \delta^{c}( q, w, 2 ) = 
 c^{q, \ell(w)}\big( q^{( k_{\ell(w)} )} \big)$. 
\end{definition}

\begin{definition}\label{reduceddefinition}
 Let $Q$ and $\Sigma$ and $\delta$ and $q_0$ and $P$ be as in the definition of a DDFA, with $\delta^{c}$ as in Definition 
 \ref{definitionchargeext}. We then define the \emph{reduced charge-extended transition function} $\overline{\delta^{c}}$ on $Q 
 \times \Sigma^{\ast}$ so that, for $q \in Q$ and for $w \in \Sigma^{\ast}$, we have that 
 $$\overline{\delta^{c}} = \delta^{c}(q, w, 1) \, \delta^{c}(q, w, 2) $$ 
 if $ \delta^{c}(q, w, 1) \in \mathbb{Q}$, 
 and, otherwise, we write 
 $\overline{\delta^{c}} $ as the formal product 
 of $ \delta^{c}(q, w, 1) $ and $\delta^{c}(q, w, 2)$, denoted as 
 $ \delta^{c}(q, w, 1) \delta^{c}(q, w, 2)$, 
 with the understanding that $q 1 = q$ and $q 0 = 0$ for $q \in Q$. 
 Simlarly, we define the \emph{reduced output function}
 $\overline{\tau}$ so that
 $$ \overline{\tau}\left( \overline{\delta^{c}}\big( q, w \big) \right) 
 = \tau\left( \delta^{c}(q, w, 1) \right) \delta^{c}(q, w, 2), $$
 with the same conventions as before concerning numerical products, and similarly for formal products. 
\end{definition}

\begin{theorem}\label{theoremmainformalize}
 Let $(Q, \Sigma, \delta, q_0, P, \Delta, \tau)$ be a DDFAO. For each $q \in Q$ and $s \in \Sigma$, let 
 $ n^{\text{current}}_{q, s} = 1$ and 
 $n_{q, t}^{\text{not current}} = 0$ for $t \in \Sigma \setminus \{ s \}$. Then 
\begin{equation}\label{20250612153PM1A}
 \overline{\delta^{c}}(q_0, w) = \delta(q_0, w) 
\end{equation}
 for an arbitrary word $w \in \Sigma^{\ast}$. 
\end{theorem}

\begin{proof}
 From Definition \ref{definitionchargeext}, we find that 
\begin{equation}\label{20250612202PM1A}
 c^{q_0, i+1}\big( q^{(k_{i})} \big) = 
 \begin{cases} 
 0 & \text{if $\delta\big( q^{(k_i)}, w_{i+1} \big) \neq q^{(k_i)} $}, \\ 
 c^{q_0, i}\big( q^{(k_{i})} \big) 
 & \text{if $\delta\big( q^{(k_i)}, w_{i+1} \big) = q^{(k_i)} $}. 
 \end{cases} 
\end{equation}
 Similarly, Definition \ref{definitionchargeext} gives us that 
\begin{equation}\label{fromdefinitioncases2}
 c^{q_0, i+1}\big( q^{(k_{i+1})} \big) = \\ 
 \begin{cases} 
 c^{q_0, i}\big( q^{(k_{i+1})} \big) + c^{q_0, i}\big( q^{(k_{i})} \big) 
 & \text{if $\delta\big( q^{(k_i)}, w_{i+1} \big) \neq q^{(k_i)} $}, \\ 
 c^{q_0, i+1}\big( q^{(k_i)} \big) & \text{if $\delta\big( q^{(k_i)}, w_{i+1} \big) = q^{(k_i)} $}, 
 \end{cases} 
\end{equation}
 and, letting $j \in \{ 1, 2, \ldots, |Q| \}$ be not equal to $k_i$ and not equal to $k_{i+1}$, we find that 
\begin{equation}\label{2025061241888888888878P8M8A}
 c^{q_0, i+1}\big( q^{(j)} \big) = c^{q_0, i}\big( q^{(j)} \big). 
 \end{equation}
 For the case whereby the empty word $w$ is empty, the desired equality in 
 \eqref{20250612153PM1A} according to the $w = \epsilon$
 cases given in the definitions for an extended transition function and for a (reduced)
 charge-extended transition function. 
 So, suppose that $w$ is nonempty. 
 We claim that 
\begin{equation}\label{claimcq0i}
 c^{q_0, i}\big( q^{(k_i)} \big) = 1
\end{equation}
 and that $c^{q_0, i}\big( q^{(j)} \big) = 0$ for $q^{(j)} \neq q^{(k_i)}$, 
 for indices $i \in \{ 0, 1, \ldots, \ell(w) \}$. 
 From Definition \ref{definitionchargeext}, 
 we find that 
 $c^{q_0, 0}(q_0) = 1$ and $c^{q_0, 0}\big( q^{(j)} \big) = 0$ 
 for $q^{(j)} \neq q_0$. Inductively, we suppose that the given claim holds for for indices $j$ not exceeding $i$. 
 Then \eqref{20250612202PM1A} reduces so that 
\begin{equation}\label{2025111000061020411PM1A}
 c^{q_0, i+1}\big( q^{(k_{i})} \big) = 
 \begin{cases} 
 0 & \text{if $\delta\big( q^{(k_i)}, w_{i+1} \big) \neq q^{(k_i)} $}, \\ 
 1 
 & \text{if $\delta\big( q^{(k_i)}, w_{i+1} \big) = q^{(k_i)} $}, 
 \end{cases} 
\end{equation}
 and \eqref{fromdefinitioncases2} reduces so that 
\begin{equation}\label{2025060122544445444P4M4A}
 c^{q_0, i+1}\big( q^{(k_{i+1})} \big) = \\ 
 \begin{cases} 
 c^{q_0, i}\big( q^{(k_{i+1})} \big) + 1 
 & \text{if $\delta\big( q^{(k_i)}, w_{i+1} \big) \neq q^{(k_i)} $}, \\ 
 c^{q_0, i+1}\big( q^{(k_i)} \big) & \text{if $\delta\big( q^{(k_i)}, w_{i+1} \big) = q^{(k_i)} $}, 
 \end{cases} 
\end{equation}
 and, moreover, since $\delta\big( q^{(k_i)}, w_{i+1} \big) \neq q^{(k_i)}$
 is equivalent to $q^{(k_{i+1})} \neq q^{(k_i)}$, we obtain from 
 \eqref{2025060122544445444P4M4A} and the inductive hypothesis that 
\begin{equation}\label{20250612411P111111111MA}
 c^{q_0, i+1}\big( q^{(k_{i+1})} \big) = \\ 
 \begin{cases} 
 1 
 & \text{if $\delta\big( q^{(k_i)}, w_{i+1} \big) \neq q^{(k_i)} $}, \\ 
 c^{q_0, i+1}\big( q^{(k_i)} \big) & \text{if $\delta\big( q^{(k_i)}, w_{i+1} \big) = q^{(k_i)} $}, 
 \end{cases} 
\end{equation}
 Moreover, from the latter case of \eqref{2025111000061020411PM1A}, 
 we find that \eqref{20250612411P111111111MA} reduces so that 
 $ c^{q_0, i+1}\big( q^{(k_{i+1})} \big) = 1$, 
 as desired. As indicated in the first case of \eqref{2025111000061020411PM1A}, we have that 
 $ c^{q_0, i+1}\big( q^{(k_{i})} \big) = 0 $ 
 if $q^{(k_i)}$ and $q^{(k_{i+1})}$ are distinct. 
 If $j \in \{ 1, 2, \ldots, |Q| \}$ is not equal to $k_i$ and is not equal to $k_{i+1}$, 
 then \eqref{2025061241888888888878P8M8A} in conjunction with the inductive hypothesis 
 give us that 
 $ c^{q_0, i+1}\big( q^{(j)} \big) = 0$, 
 as desired. The definition in \eqref{setdeltacqw} 
 together with \eqref{claimcq0i} then give us that 
\begin{equation}\label{pairwith1}
 \delta^{c}( q_0, w ) 
 = \left( q^{( k_{\ell(w)} )}, 
 1 \right), 
\end{equation}
 so that \eqref{pairwith1} and Definition \ref{reduceddefinition} together give us that 
\begin{equation}\label{pairreducedtogether}
 \overline{\delta^{c}}( q_0, w ) 
 = q^{( k_{\ell(w)} )}, 
\end{equation}
 and the right-hand side of \eqref{pairreducedtogether} is equal to 
 $\delta(q_0, w)$, according to \eqref{deltanested}. 
\end{proof}

\end{document}
